% =====================================================================
% Paper 07 — Per-generation quantum consistency: IDEA / MECHANISM NOTE
%   on Ward + Einstein--Hilbert + Witten + mod-2 framework
%   (NOT a current theorem paper; external use FORBIDDEN at present
%    canonical tier per Math314-AddC audit; full theorem-level rewrite
%    tracked under Q-2026-05-02-Pillar7-Three-Generation-Rewrite)
% Status: [DRAFT] [REWRITE_REQUIRED] [EXTERNAL_USE_FORBIDDEN]
% AUDIT-FLAG: AUDIT-2026-05-02-Wave7-Aux-Epoch-Overclaim
%   (Math314-AddC extension covering Wave 4/5 papers Paper-06..08).
% Rev 4 (2026-05-06): full body rewrite as mechanism note.
%   Specific corrections applied throughout the body:
%   (a) Introduction wording "the unique solution is three generations"
%       removed; replaced by an honest framing: TECT identifies the
%       SM three-generation chiral spectrum as anomaly-cancellation-
%       consistent within SO(10), but does NOT prove that exactly
%       three generations are forced.
%   (b) Theorem statement no longer asserts "exactly three generations
%       are required by quantum consistency"; it is recast as a
%       four-condition consistency-conjunction conjecture (Math47 Ward
%       identity + Math48 nonlinear EH coupling + Witten SU(2) global
%       + Math49c mod-2 spectral flow), with the generation-count
%       question explicitly identified as canonical Pillar-7 T1 OPEN.
%   (c) Proof "Step 2 (Nonlinear EH coupling)" removed: the formula
%       A_grav = sum_f Q_f^3 (sum of electric charges cubed) is NOT a
%       gravitational anomaly coefficient. The actual mixed
%       gravitational anomaly is Tr(U(1)_Y) (one power of Y, not three
%       powers of Q), which vanishes per generation in the SM
%       independent of N_g (see Paper-07-ext for the canonical
%       cancellation). The Math48 nonlinear EH content is recast as a
%       per-generation Ward-current preservation under coupling to a
%       fluctuating background metric, NOT a forcing condition on N_g.
%   (d) "Step 3 (Witten SU(2) anomaly)" rewritten precisely. Witten's
%       theorem is a Z_2 statement on the chiral SU(2)_L doublet
%       count modulo 2; it does NOT distinguish N_g = 1 from N_g = 3,
%       since both give an even total doublet count (4 doublets per
%       generation x N_g). The mis-attribution "the coefficient is
%       zero iff N_g = 3" is corrected.
%   (e) "Step 4 (Mod-2 spectral flow)" rewritten precisely. The
%       parity formula (-1)^{N_g(N_g-1)/2} appearing in the earlier
%       draft is non-standard and not tied to any TECT canonical
%       construction. It is replaced by a faithful summary of
%       Math49c-v3's mod-2 spectral-flow construction: the chiral
%       Dirac operator on the BCC shell admits an integer-valued
%       Atiyah-Singer index whose mod-2 reduction must be even for
%       quantization consistency; the question of whether this index
%       picks out N_g = 3 is part of Math49c-v3 sub-task and remains
%       Pillar-7 T1 OPEN.
%   (f) "Closure" paragraph recast: the four conditions are necessary
%       per-generation consistency conditions; together they are not
%       known to be sufficient to FORCE N_g = 3.
%   (g) Discussion "result is PROVED (T7)" REMOVED. The honest tier
%       is T2 CONJECTURE for the three-generation-forcing claim, with
%       four-condition mechanism summary as the canonical content of
%       this paper.
% Compile: bash _shared/build-paper.sh Paper-07-Pillar7-Quantum-Consistency/Paper-07
% Canonical archive: Math47, Math48, Math49b-v3, Math49c-v3,
%                    Math314-AddC, Paper-07-ext (anomaly-cancellation
%                    auxiliary)
% =====================================================================
% --- TECT canonical preamble (see _shared/tect-preamble.tex) ----------
% =====================================================================
% TECT canonical LaTeX preamble (revtex4-2 / single-column / theorem-safe)
% =====================================================================
% Maintainer: Jusang Lee (jtkor@outlook.com)
% Binding from: 2026-05-06
% Purpose: a single source-of-truth preamble for every TECT paper
% (Paper-00 .. Paper-10 + Auxiliary notes). Designed to (a) eliminate
% font-corruption on pdflatex (T1 fontenc + lmodern), (b) make
% theorem-heavy mathematical-physics content render cleanly in a
% single-column layout (PRD class option, NOT PRL two-column), (c)
% provide consistent theorem numbering across all TECT papers via
% amsthm with a shared counter set, (d) make hyperref + cleveref
% cross-references resolve cleanly with the standard two-pass
% pdflatex compile.
%
% Usage in each Paper-NN.tex:
%   % =====================================================================
% TECT canonical LaTeX preamble (revtex4-2 / single-column / theorem-safe)
% =====================================================================
% Maintainer: Jusang Lee (jtkor@outlook.com)
% Binding from: 2026-05-06
% Purpose: a single source-of-truth preamble for every TECT paper
% (Paper-00 .. Paper-10 + Auxiliary notes). Designed to (a) eliminate
% font-corruption on pdflatex (T1 fontenc + lmodern), (b) make
% theorem-heavy mathematical-physics content render cleanly in a
% single-column layout (PRD class option, NOT PRL two-column), (c)
% provide consistent theorem numbering across all TECT papers via
% amsthm with a shared counter set, (d) make hyperref + cleveref
% cross-references resolve cleanly with the standard two-pass
% pdflatex compile.
%
% Usage in each Paper-NN.tex:
%   % =====================================================================
% TECT canonical LaTeX preamble (revtex4-2 / single-column / theorem-safe)
% =====================================================================
% Maintainer: Jusang Lee (jtkor@outlook.com)
% Binding from: 2026-05-06
% Purpose: a single source-of-truth preamble for every TECT paper
% (Paper-00 .. Paper-10 + Auxiliary notes). Designed to (a) eliminate
% font-corruption on pdflatex (T1 fontenc + lmodern), (b) make
% theorem-heavy mathematical-physics content render cleanly in a
% single-column layout (PRD class option, NOT PRL two-column), (c)
% provide consistent theorem numbering across all TECT papers via
% amsthm with a shared counter set, (d) make hyperref + cleveref
% cross-references resolve cleanly with the standard two-pass
% pdflatex compile.
%
% Usage in each Paper-NN.tex:
%   \input{../_shared/tect-preamble.tex}
%   \begin{document}
%   ...
%   \end{document}
%
% Build (every paper): `pdflatex Paper-NN && pdflatex Paper-NN`
% (the second pass resolves \ref{} / \cref{} forward references).
% A bash/PowerShell wrapper that automates this is at
% Docs/papers/papers/_shared/build-paper.sh / build-paper.ps1.
% =====================================================================

% ---- Document class --------------------------------------------------
% PRD class option of revtex4-2 with [reprint,onecolumn] gives the same
% APS heading / by-line / abstract style as PRL but in a wider single
% column that handles long display equations and theorem environments
% without column-break artefacts. nofootinbib keeps citations out of
% footnotes. The 11pt option restores standard font metrics; with T1
% fontenc + lmodern this gives non-bitmap glyphs.
\documentclass[%
  aps,prd,reprint,onecolumn,nofootinbib,11pt,%
  amsmath,amssymb,floatfix,longbibliography%
]{revtex4-2}

% ---- Encoding & fonts ------------------------------------------------
\usepackage[T1]{fontenc}        % proper 8-bit font encoding (no font corruption)
\usepackage[utf8]{inputenc}     % allow UTF-8 source
\usepackage{lmodern}            % scalable Latin Modern (replaces bitmap CM)
\usepackage{textcomp}           % extra text symbols
\usepackage{microtype}          % subtle kerning improvements

% ---- UTF-8 → LaTeX character mappings --------------------------------
% Maps the Unicode characters that occur in TECT body text (legacy from
% the chat-content auto-archival rule, CLAUDE.md §4) to LaTeX-safe
% commands. Without these declarations, T1+inputenc rejects the chars
% with "Unicode character X not set up for use with LaTeX".
% Adding declarations centrally (here) is preferable to per-file edits.
\DeclareUnicodeCharacter{00A7}{\S}              % §  (section sign)
\DeclareUnicodeCharacter{00B2}{\ensuremath{^{2}}}        % ²
\DeclareUnicodeCharacter{00B3}{\ensuremath{^{3}}}        % ³
\DeclareUnicodeCharacter{00B5}{\ensuremath{\mu}}         % µ (micro)
\DeclareUnicodeCharacter{00B7}{\ensuremath{\cdot}}       % · (middle dot)
\DeclareUnicodeCharacter{00D7}{\ensuremath{\times}}      % ×
\DeclareUnicodeCharacter{00E4}{\"a}             % ä
\DeclareUnicodeCharacter{00E9}{\'e}             % é
\DeclareUnicodeCharacter{00F6}{\"o}             % ö
\DeclareUnicodeCharacter{00FC}{\"u}             % ü
\DeclareUnicodeCharacter{010C}{\v{C}}           % Č
\DeclareUnicodeCharacter{010D}{\v{c}}           % č
\DeclareUnicodeCharacter{0394}{\ensuremath{\Delta}}      % Δ
\DeclareUnicodeCharacter{03A3}{\ensuremath{\Sigma}}      % Σ
\DeclareUnicodeCharacter{03A9}{\ensuremath{\Omega}}      % Ω
\DeclareUnicodeCharacter{03B1}{\ensuremath{\alpha}}      % α
\DeclareUnicodeCharacter{03B2}{\ensuremath{\beta}}       % β
\DeclareUnicodeCharacter{03B3}{\ensuremath{\gamma}}      % γ
\DeclareUnicodeCharacter{03B4}{\ensuremath{\delta}}      % δ
\DeclareUnicodeCharacter{03B5}{\ensuremath{\varepsilon}} % ε
\DeclareUnicodeCharacter{03BA}{\ensuremath{\kappa}}      % κ
\DeclareUnicodeCharacter{03BB}{\ensuremath{\lambda}}     % λ
\DeclareUnicodeCharacter{03BC}{\ensuremath{\mu}}         % μ
\DeclareUnicodeCharacter{03BD}{\ensuremath{\nu}}         % ν
\DeclareUnicodeCharacter{03C0}{\ensuremath{\pi}}         % π
\DeclareUnicodeCharacter{03C1}{\ensuremath{\rho}}        % ρ
\DeclareUnicodeCharacter{03C3}{\ensuremath{\sigma}}      % σ
\DeclareUnicodeCharacter{03C4}{\ensuremath{\tau}}        % τ
\DeclareUnicodeCharacter{03C6}{\ensuremath{\varphi}}     % φ
\DeclareUnicodeCharacter{03C7}{\ensuremath{\chi}}        % χ
\DeclareUnicodeCharacter{03C8}{\ensuremath{\psi}}        % ψ
\DeclareUnicodeCharacter{03C9}{\ensuremath{\omega}}      % ω
\DeclareUnicodeCharacter{2013}{--}              % – (en dash)
\DeclareUnicodeCharacter{2014}{---}             % — (em dash)
\DeclareUnicodeCharacter{2018}{`}               % ‘
\DeclareUnicodeCharacter{2019}{'}               % ’
\DeclareUnicodeCharacter{201C}{``}              % “
\DeclareUnicodeCharacter{201D}{''}              % ”
\DeclareUnicodeCharacter{2070}{\ensuremath{^{0}}}        % ⁰
\DeclareUnicodeCharacter{2071}{\ensuremath{^{i}}}        % ⁱ
\DeclareUnicodeCharacter{2122}{\textsuperscript{TM}}     % ™
\DeclareUnicodeCharacter{2126}{\ensuremath{\Omega}}      % Ω (ohm sign)
\DeclareUnicodeCharacter{210F}{\ensuremath{\hbar}}       % ℏ
\DeclareUnicodeCharacter{2115}{\ensuremath{\mathbb{N}}}  % ℕ
\DeclareUnicodeCharacter{2102}{\ensuremath{\mathbb{C}}}  % ℂ
\DeclareUnicodeCharacter{211D}{\ensuremath{\mathbb{R}}}  % ℝ
\DeclareUnicodeCharacter{2124}{\ensuremath{\mathbb{Z}}}  % ℤ
\DeclareUnicodeCharacter{211A}{\ensuremath{\mathbb{Q}}}  % ℚ
\DeclareUnicodeCharacter{2192}{\ensuremath{\to}}         % →
\DeclareUnicodeCharacter{21D2}{\ensuremath{\Rightarrow}} % ⇒
\DeclareUnicodeCharacter{2200}{\ensuremath{\forall}}     % ∀
\DeclareUnicodeCharacter{2203}{\ensuremath{\exists}}     % ∃
\DeclareUnicodeCharacter{2208}{\ensuremath{\in}}         % ∈
\DeclareUnicodeCharacter{2209}{\ensuremath{\notin}}      % ∉
\DeclareUnicodeCharacter{2227}{\ensuremath{\wedge}}      % ∧
\DeclareUnicodeCharacter{2228}{\ensuremath{\vee}}        % ∨
\DeclareUnicodeCharacter{2229}{\ensuremath{\cap}}        % ∩
\DeclareUnicodeCharacter{222A}{\ensuremath{\cup}}        % ∪
\DeclareUnicodeCharacter{2245}{\ensuremath{\cong}}       % ≅
\DeclareUnicodeCharacter{2248}{\ensuremath{\approx}}     % ≈
\DeclareUnicodeCharacter{2260}{\ensuremath{\neq}}        % ≠
\DeclareUnicodeCharacter{2264}{\ensuremath{\leq}}        % ≤
\DeclareUnicodeCharacter{2265}{\ensuremath{\geq}}        % ≥
\DeclareUnicodeCharacter{22C5}{\ensuremath{\cdot}}       % ⋅
\DeclareUnicodeCharacter{2295}{\ensuremath{\oplus}}      % ⊕
\DeclareUnicodeCharacter{2297}{\ensuremath{\otimes}}     % ⊗

% ---- Mathematics & symbols ------------------------------------------
\usepackage{amsmath,amssymb,bm,mathrsfs,mathtools}
\usepackage{amsthm}             % theorem environments (load BEFORE hyperref)

% ---- Theorem environments (shared counter set, TECT canonical) -------
% All theorems / lemmas / propositions / corollaries / remarks share a
% single counter so cross-references read consistently across a paper.
% Definitions get a separate counter (definitions are numbered in their
% own sequence by editorial convention).
\newtheorem{theorem}{Theorem}
\newtheorem{lemma}[theorem]{Lemma}
\newtheorem{proposition}[theorem]{Proposition}
\newtheorem{corollary}[theorem]{Corollary}

\theoremstyle{remark}
\newtheorem{remark}[theorem]{Remark}
\newtheorem{example}[theorem]{Example}

\theoremstyle{definition}
\newtheorem{definition}[theorem]{Definition}
\newtheorem{conjecture}[theorem]{Conjecture}
\newtheorem{hypothesis}[theorem]{Hypothesis}

% ---- Tables / floats / graphics --------------------------------------
\usepackage{booktabs}           % \toprule / \midrule / \bottomrule
\usepackage{array}
\usepackage{graphicx}
\usepackage{enumitem}           % \begin{itemize}[leftmargin=...] options
\usepackage{hyphenat}           % \hyp{} for breakable hyphens in \texttt
\usepackage{seqsplit}           % \seqsplit{} for long unbreakable strings
\sloppy                         % allow stretchier inter-word spacing to
                                % reduce overfull \hbox warnings on long URLs
                                % and long \texttt tokens (paragraph-level).
\emergencystretch=3em           % last-resort hbox stretch budget

% ---- Cross-references (hyperref last among the standard packages,
%      cleveref AFTER hyperref) ----------------------------------------
\usepackage[%
  colorlinks=true,%
  linkcolor=blue,%
  citecolor=blue,%
  urlcolor=blue,%
  breaklinks=true,%
  bookmarks=true,%
  bookmarksnumbered=true%
]{hyperref}
\usepackage[capitalise,nameinlink]{cleveref}

% ---- TECT-canonical author block macros -----------------------------
% (defined here so every paper has the same author / affiliation style)
\newcommand{\tectauthor}{%
  \author{Jusang Lee}%
  \email{jtkor@outlook.com}%
  \affiliation{Independent researcher, Republic of Korea}%
}

% ---- TECT-canonical archive footer ----------------------------------
\newcommand{\tectarchivefooter}{%
  \par\medskip
  \noindent\small\textit{Canonical archive}:
  \url{https://github.com/TECT-OpenPhysics/TECT}.\par%
}

% ---- TECT TODO / AUDIT-FLAG / HONEST-SCOPE inline marks --------------
% Used in drafts to highlight pending audit items without disturbing
% the body text style. Suppress via \tectsuppressmarks (define empty).
\providecommand{\tectauditflag}[1]{\textcolor{red}{\textbf{[AUDIT-FLAG: #1]}}}
\providecommand{\tecttodo}[1]{\textcolor{orange}{\textbf{[TODO: #1]}}}
\providecommand{\tecthonestscope}[1]{\textit{Honest scope: #1.}}

% =====================================================================
% End of TECT canonical preamble
% =====================================================================

%   \begin{document}
%   ...
%   \end{document}
%
% Build (every paper): `pdflatex Paper-NN && pdflatex Paper-NN`
% (the second pass resolves \ref{} / \cref{} forward references).
% A bash/PowerShell wrapper that automates this is at
% Docs/papers/papers/_shared/build-paper.sh / build-paper.ps1.
% =====================================================================

% ---- Document class --------------------------------------------------
% PRD class option of revtex4-2 with [reprint,onecolumn] gives the same
% APS heading / by-line / abstract style as PRL but in a wider single
% column that handles long display equations and theorem environments
% without column-break artefacts. nofootinbib keeps citations out of
% footnotes. The 11pt option restores standard font metrics; with T1
% fontenc + lmodern this gives non-bitmap glyphs.
\documentclass[%
  aps,prd,reprint,onecolumn,nofootinbib,11pt,%
  amsmath,amssymb,floatfix,longbibliography%
]{revtex4-2}

% ---- Encoding & fonts ------------------------------------------------
\usepackage[T1]{fontenc}        % proper 8-bit font encoding (no font corruption)
\usepackage[utf8]{inputenc}     % allow UTF-8 source
\usepackage{lmodern}            % scalable Latin Modern (replaces bitmap CM)
\usepackage{textcomp}           % extra text symbols
\usepackage{microtype}          % subtle kerning improvements

% ---- UTF-8 → LaTeX character mappings --------------------------------
% Maps the Unicode characters that occur in TECT body text (legacy from
% the chat-content auto-archival rule, CLAUDE.md §4) to LaTeX-safe
% commands. Without these declarations, T1+inputenc rejects the chars
% with "Unicode character X not set up for use with LaTeX".
% Adding declarations centrally (here) is preferable to per-file edits.
\DeclareUnicodeCharacter{00A7}{\S}              % §  (section sign)
\DeclareUnicodeCharacter{00B2}{\ensuremath{^{2}}}        % ²
\DeclareUnicodeCharacter{00B3}{\ensuremath{^{3}}}        % ³
\DeclareUnicodeCharacter{00B5}{\ensuremath{\mu}}         % µ (micro)
\DeclareUnicodeCharacter{00B7}{\ensuremath{\cdot}}       % · (middle dot)
\DeclareUnicodeCharacter{00D7}{\ensuremath{\times}}      % ×
\DeclareUnicodeCharacter{00E4}{\"a}             % ä
\DeclareUnicodeCharacter{00E9}{\'e}             % é
\DeclareUnicodeCharacter{00F6}{\"o}             % ö
\DeclareUnicodeCharacter{00FC}{\"u}             % ü
\DeclareUnicodeCharacter{010C}{\v{C}}           % Č
\DeclareUnicodeCharacter{010D}{\v{c}}           % č
\DeclareUnicodeCharacter{0394}{\ensuremath{\Delta}}      % Δ
\DeclareUnicodeCharacter{03A3}{\ensuremath{\Sigma}}      % Σ
\DeclareUnicodeCharacter{03A9}{\ensuremath{\Omega}}      % Ω
\DeclareUnicodeCharacter{03B1}{\ensuremath{\alpha}}      % α
\DeclareUnicodeCharacter{03B2}{\ensuremath{\beta}}       % β
\DeclareUnicodeCharacter{03B3}{\ensuremath{\gamma}}      % γ
\DeclareUnicodeCharacter{03B4}{\ensuremath{\delta}}      % δ
\DeclareUnicodeCharacter{03B5}{\ensuremath{\varepsilon}} % ε
\DeclareUnicodeCharacter{03BA}{\ensuremath{\kappa}}      % κ
\DeclareUnicodeCharacter{03BB}{\ensuremath{\lambda}}     % λ
\DeclareUnicodeCharacter{03BC}{\ensuremath{\mu}}         % μ
\DeclareUnicodeCharacter{03BD}{\ensuremath{\nu}}         % ν
\DeclareUnicodeCharacter{03C0}{\ensuremath{\pi}}         % π
\DeclareUnicodeCharacter{03C1}{\ensuremath{\rho}}        % ρ
\DeclareUnicodeCharacter{03C3}{\ensuremath{\sigma}}      % σ
\DeclareUnicodeCharacter{03C4}{\ensuremath{\tau}}        % τ
\DeclareUnicodeCharacter{03C6}{\ensuremath{\varphi}}     % φ
\DeclareUnicodeCharacter{03C7}{\ensuremath{\chi}}        % χ
\DeclareUnicodeCharacter{03C8}{\ensuremath{\psi}}        % ψ
\DeclareUnicodeCharacter{03C9}{\ensuremath{\omega}}      % ω
\DeclareUnicodeCharacter{2013}{--}              % – (en dash)
\DeclareUnicodeCharacter{2014}{---}             % — (em dash)
\DeclareUnicodeCharacter{2018}{`}               % ‘
\DeclareUnicodeCharacter{2019}{'}               % ’
\DeclareUnicodeCharacter{201C}{``}              % “
\DeclareUnicodeCharacter{201D}{''}              % ”
\DeclareUnicodeCharacter{2070}{\ensuremath{^{0}}}        % ⁰
\DeclareUnicodeCharacter{2071}{\ensuremath{^{i}}}        % ⁱ
\DeclareUnicodeCharacter{2122}{\textsuperscript{TM}}     % ™
\DeclareUnicodeCharacter{2126}{\ensuremath{\Omega}}      % Ω (ohm sign)
\DeclareUnicodeCharacter{210F}{\ensuremath{\hbar}}       % ℏ
\DeclareUnicodeCharacter{2115}{\ensuremath{\mathbb{N}}}  % ℕ
\DeclareUnicodeCharacter{2102}{\ensuremath{\mathbb{C}}}  % ℂ
\DeclareUnicodeCharacter{211D}{\ensuremath{\mathbb{R}}}  % ℝ
\DeclareUnicodeCharacter{2124}{\ensuremath{\mathbb{Z}}}  % ℤ
\DeclareUnicodeCharacter{211A}{\ensuremath{\mathbb{Q}}}  % ℚ
\DeclareUnicodeCharacter{2192}{\ensuremath{\to}}         % →
\DeclareUnicodeCharacter{21D2}{\ensuremath{\Rightarrow}} % ⇒
\DeclareUnicodeCharacter{2200}{\ensuremath{\forall}}     % ∀
\DeclareUnicodeCharacter{2203}{\ensuremath{\exists}}     % ∃
\DeclareUnicodeCharacter{2208}{\ensuremath{\in}}         % ∈
\DeclareUnicodeCharacter{2209}{\ensuremath{\notin}}      % ∉
\DeclareUnicodeCharacter{2227}{\ensuremath{\wedge}}      % ∧
\DeclareUnicodeCharacter{2228}{\ensuremath{\vee}}        % ∨
\DeclareUnicodeCharacter{2229}{\ensuremath{\cap}}        % ∩
\DeclareUnicodeCharacter{222A}{\ensuremath{\cup}}        % ∪
\DeclareUnicodeCharacter{2245}{\ensuremath{\cong}}       % ≅
\DeclareUnicodeCharacter{2248}{\ensuremath{\approx}}     % ≈
\DeclareUnicodeCharacter{2260}{\ensuremath{\neq}}        % ≠
\DeclareUnicodeCharacter{2264}{\ensuremath{\leq}}        % ≤
\DeclareUnicodeCharacter{2265}{\ensuremath{\geq}}        % ≥
\DeclareUnicodeCharacter{22C5}{\ensuremath{\cdot}}       % ⋅
\DeclareUnicodeCharacter{2295}{\ensuremath{\oplus}}      % ⊕
\DeclareUnicodeCharacter{2297}{\ensuremath{\otimes}}     % ⊗

% ---- Mathematics & symbols ------------------------------------------
\usepackage{amsmath,amssymb,bm,mathrsfs,mathtools}
\usepackage{amsthm}             % theorem environments (load BEFORE hyperref)

% ---- Theorem environments (shared counter set, TECT canonical) -------
% All theorems / lemmas / propositions / corollaries / remarks share a
% single counter so cross-references read consistently across a paper.
% Definitions get a separate counter (definitions are numbered in their
% own sequence by editorial convention).
\newtheorem{theorem}{Theorem}
\newtheorem{lemma}[theorem]{Lemma}
\newtheorem{proposition}[theorem]{Proposition}
\newtheorem{corollary}[theorem]{Corollary}

\theoremstyle{remark}
\newtheorem{remark}[theorem]{Remark}
\newtheorem{example}[theorem]{Example}

\theoremstyle{definition}
\newtheorem{definition}[theorem]{Definition}
\newtheorem{conjecture}[theorem]{Conjecture}
\newtheorem{hypothesis}[theorem]{Hypothesis}

% ---- Tables / floats / graphics --------------------------------------
\usepackage{booktabs}           % \toprule / \midrule / \bottomrule
\usepackage{array}
\usepackage{graphicx}
\usepackage{enumitem}           % \begin{itemize}[leftmargin=...] options
\usepackage{hyphenat}           % \hyp{} for breakable hyphens in \texttt
\usepackage{seqsplit}           % \seqsplit{} for long unbreakable strings
\sloppy                         % allow stretchier inter-word spacing to
                                % reduce overfull \hbox warnings on long URLs
                                % and long \texttt tokens (paragraph-level).
\emergencystretch=3em           % last-resort hbox stretch budget

% ---- Cross-references (hyperref last among the standard packages,
%      cleveref AFTER hyperref) ----------------------------------------
\usepackage[%
  colorlinks=true,%
  linkcolor=blue,%
  citecolor=blue,%
  urlcolor=blue,%
  breaklinks=true,%
  bookmarks=true,%
  bookmarksnumbered=true%
]{hyperref}
\usepackage[capitalise,nameinlink]{cleveref}

% ---- TECT-canonical author block macros -----------------------------
% (defined here so every paper has the same author / affiliation style)
\newcommand{\tectauthor}{%
  \author{Jusang Lee}%
  \email{jtkor@outlook.com}%
  \affiliation{Independent researcher, Republic of Korea}%
}

% ---- TECT-canonical archive footer ----------------------------------
\newcommand{\tectarchivefooter}{%
  \par\medskip
  \noindent\small\textit{Canonical archive}:
  \url{https://github.com/TECT-OpenPhysics/TECT}.\par%
}

% ---- TECT TODO / AUDIT-FLAG / HONEST-SCOPE inline marks --------------
% Used in drafts to highlight pending audit items without disturbing
% the body text style. Suppress via \tectsuppressmarks (define empty).
\providecommand{\tectauditflag}[1]{\textcolor{red}{\textbf{[AUDIT-FLAG: #1]}}}
\providecommand{\tecttodo}[1]{\textcolor{orange}{\textbf{[TODO: #1]}}}
\providecommand{\tecthonestscope}[1]{\textit{Honest scope: #1.}}

% =====================================================================
% End of TECT canonical preamble
% =====================================================================

%   \begin{document}
%   ...
%   \end{document}
%
% Build (every paper): `pdflatex Paper-NN && pdflatex Paper-NN`
% (the second pass resolves \ref{} / \cref{} forward references).
% A bash/PowerShell wrapper that automates this is at
% Docs/papers/papers/_shared/build-paper.sh / build-paper.ps1.
% =====================================================================

% ---- Document class --------------------------------------------------
% PRD class option of revtex4-2 with [reprint,onecolumn] gives the same
% APS heading / by-line / abstract style as PRL but in a wider single
% column that handles long display equations and theorem environments
% without column-break artefacts. nofootinbib keeps citations out of
% footnotes. The 11pt option restores standard font metrics; with T1
% fontenc + lmodern this gives non-bitmap glyphs.
\documentclass[%
  aps,prd,reprint,onecolumn,nofootinbib,11pt,%
  amsmath,amssymb,floatfix,longbibliography%
]{revtex4-2}

% ---- Encoding & fonts ------------------------------------------------
\usepackage[T1]{fontenc}        % proper 8-bit font encoding (no font corruption)
\usepackage[utf8]{inputenc}     % allow UTF-8 source
\usepackage{lmodern}            % scalable Latin Modern (replaces bitmap CM)
\usepackage{textcomp}           % extra text symbols
\usepackage{microtype}          % subtle kerning improvements

% ---- UTF-8 → LaTeX character mappings --------------------------------
% Maps the Unicode characters that occur in TECT body text (legacy from
% the chat-content auto-archival rule, CLAUDE.md §4) to LaTeX-safe
% commands. Without these declarations, T1+inputenc rejects the chars
% with "Unicode character X not set up for use with LaTeX".
% Adding declarations centrally (here) is preferable to per-file edits.
\DeclareUnicodeCharacter{00A7}{\S}              % §  (section sign)
\DeclareUnicodeCharacter{00B2}{\ensuremath{^{2}}}        % ²
\DeclareUnicodeCharacter{00B3}{\ensuremath{^{3}}}        % ³
\DeclareUnicodeCharacter{00B5}{\ensuremath{\mu}}         % µ (micro)
\DeclareUnicodeCharacter{00B7}{\ensuremath{\cdot}}       % · (middle dot)
\DeclareUnicodeCharacter{00D7}{\ensuremath{\times}}      % ×
\DeclareUnicodeCharacter{00E4}{\"a}             % ä
\DeclareUnicodeCharacter{00E9}{\'e}             % é
\DeclareUnicodeCharacter{00F6}{\"o}             % ö
\DeclareUnicodeCharacter{00FC}{\"u}             % ü
\DeclareUnicodeCharacter{010C}{\v{C}}           % Č
\DeclareUnicodeCharacter{010D}{\v{c}}           % č
\DeclareUnicodeCharacter{0394}{\ensuremath{\Delta}}      % Δ
\DeclareUnicodeCharacter{03A3}{\ensuremath{\Sigma}}      % Σ
\DeclareUnicodeCharacter{03A9}{\ensuremath{\Omega}}      % Ω
\DeclareUnicodeCharacter{03B1}{\ensuremath{\alpha}}      % α
\DeclareUnicodeCharacter{03B2}{\ensuremath{\beta}}       % β
\DeclareUnicodeCharacter{03B3}{\ensuremath{\gamma}}      % γ
\DeclareUnicodeCharacter{03B4}{\ensuremath{\delta}}      % δ
\DeclareUnicodeCharacter{03B5}{\ensuremath{\varepsilon}} % ε
\DeclareUnicodeCharacter{03BA}{\ensuremath{\kappa}}      % κ
\DeclareUnicodeCharacter{03BB}{\ensuremath{\lambda}}     % λ
\DeclareUnicodeCharacter{03BC}{\ensuremath{\mu}}         % μ
\DeclareUnicodeCharacter{03BD}{\ensuremath{\nu}}         % ν
\DeclareUnicodeCharacter{03C0}{\ensuremath{\pi}}         % π
\DeclareUnicodeCharacter{03C1}{\ensuremath{\rho}}        % ρ
\DeclareUnicodeCharacter{03C3}{\ensuremath{\sigma}}      % σ
\DeclareUnicodeCharacter{03C4}{\ensuremath{\tau}}        % τ
\DeclareUnicodeCharacter{03C6}{\ensuremath{\varphi}}     % φ
\DeclareUnicodeCharacter{03C7}{\ensuremath{\chi}}        % χ
\DeclareUnicodeCharacter{03C8}{\ensuremath{\psi}}        % ψ
\DeclareUnicodeCharacter{03C9}{\ensuremath{\omega}}      % ω
\DeclareUnicodeCharacter{2013}{--}              % – (en dash)
\DeclareUnicodeCharacter{2014}{---}             % — (em dash)
\DeclareUnicodeCharacter{2018}{`}               % ‘
\DeclareUnicodeCharacter{2019}{'}               % ’
\DeclareUnicodeCharacter{201C}{``}              % “
\DeclareUnicodeCharacter{201D}{''}              % ”
\DeclareUnicodeCharacter{2070}{\ensuremath{^{0}}}        % ⁰
\DeclareUnicodeCharacter{2071}{\ensuremath{^{i}}}        % ⁱ
\DeclareUnicodeCharacter{2122}{\textsuperscript{TM}}     % ™
\DeclareUnicodeCharacter{2126}{\ensuremath{\Omega}}      % Ω (ohm sign)
\DeclareUnicodeCharacter{210F}{\ensuremath{\hbar}}       % ℏ
\DeclareUnicodeCharacter{2115}{\ensuremath{\mathbb{N}}}  % ℕ
\DeclareUnicodeCharacter{2102}{\ensuremath{\mathbb{C}}}  % ℂ
\DeclareUnicodeCharacter{211D}{\ensuremath{\mathbb{R}}}  % ℝ
\DeclareUnicodeCharacter{2124}{\ensuremath{\mathbb{Z}}}  % ℤ
\DeclareUnicodeCharacter{211A}{\ensuremath{\mathbb{Q}}}  % ℚ
\DeclareUnicodeCharacter{2192}{\ensuremath{\to}}         % →
\DeclareUnicodeCharacter{21D2}{\ensuremath{\Rightarrow}} % ⇒
\DeclareUnicodeCharacter{2200}{\ensuremath{\forall}}     % ∀
\DeclareUnicodeCharacter{2203}{\ensuremath{\exists}}     % ∃
\DeclareUnicodeCharacter{2208}{\ensuremath{\in}}         % ∈
\DeclareUnicodeCharacter{2209}{\ensuremath{\notin}}      % ∉
\DeclareUnicodeCharacter{2227}{\ensuremath{\wedge}}      % ∧
\DeclareUnicodeCharacter{2228}{\ensuremath{\vee}}        % ∨
\DeclareUnicodeCharacter{2229}{\ensuremath{\cap}}        % ∩
\DeclareUnicodeCharacter{222A}{\ensuremath{\cup}}        % ∪
\DeclareUnicodeCharacter{2245}{\ensuremath{\cong}}       % ≅
\DeclareUnicodeCharacter{2248}{\ensuremath{\approx}}     % ≈
\DeclareUnicodeCharacter{2260}{\ensuremath{\neq}}        % ≠
\DeclareUnicodeCharacter{2264}{\ensuremath{\leq}}        % ≤
\DeclareUnicodeCharacter{2265}{\ensuremath{\geq}}        % ≥
\DeclareUnicodeCharacter{22C5}{\ensuremath{\cdot}}       % ⋅
\DeclareUnicodeCharacter{2295}{\ensuremath{\oplus}}      % ⊕
\DeclareUnicodeCharacter{2297}{\ensuremath{\otimes}}     % ⊗

% ---- Mathematics & symbols ------------------------------------------
\usepackage{amsmath,amssymb,bm,mathrsfs,mathtools}
\usepackage{amsthm}             % theorem environments (load BEFORE hyperref)

% ---- Theorem environments (shared counter set, TECT canonical) -------
% All theorems / lemmas / propositions / corollaries / remarks share a
% single counter so cross-references read consistently across a paper.
% Definitions get a separate counter (definitions are numbered in their
% own sequence by editorial convention).
\newtheorem{theorem}{Theorem}
\newtheorem{lemma}[theorem]{Lemma}
\newtheorem{proposition}[theorem]{Proposition}
\newtheorem{corollary}[theorem]{Corollary}

\theoremstyle{remark}
\newtheorem{remark}[theorem]{Remark}
\newtheorem{example}[theorem]{Example}

\theoremstyle{definition}
\newtheorem{definition}[theorem]{Definition}
\newtheorem{conjecture}[theorem]{Conjecture}
\newtheorem{hypothesis}[theorem]{Hypothesis}

% ---- Tables / floats / graphics --------------------------------------
\usepackage{booktabs}           % \toprule / \midrule / \bottomrule
\usepackage{array}
\usepackage{graphicx}
\usepackage{enumitem}           % \begin{itemize}[leftmargin=...] options
\usepackage{hyphenat}           % \hyp{} for breakable hyphens in \texttt
\usepackage{seqsplit}           % \seqsplit{} for long unbreakable strings
\sloppy                         % allow stretchier inter-word spacing to
                                % reduce overfull \hbox warnings on long URLs
                                % and long \texttt tokens (paragraph-level).
\emergencystretch=3em           % last-resort hbox stretch budget

% ---- Cross-references (hyperref last among the standard packages,
%      cleveref AFTER hyperref) ----------------------------------------
\usepackage[%
  colorlinks=true,%
  linkcolor=blue,%
  citecolor=blue,%
  urlcolor=blue,%
  breaklinks=true,%
  bookmarks=true,%
  bookmarksnumbered=true%
]{hyperref}
\usepackage[capitalise,nameinlink]{cleveref}

% ---- TECT-canonical author block macros -----------------------------
% (defined here so every paper has the same author / affiliation style)
\newcommand{\tectauthor}{%
  \author{Jusang Lee}%
  \email{jtkor@outlook.com}%
  \affiliation{Independent researcher, Republic of Korea}%
}

% ---- TECT-canonical archive footer ----------------------------------
\newcommand{\tectarchivefooter}{%
  \par\medskip
  \noindent\small\textit{Canonical archive}:
  \url{https://github.com/TECT-OpenPhysics/TECT}.\par%
}

% ---- TECT TODO / AUDIT-FLAG / HONEST-SCOPE inline marks --------------
% Used in drafts to highlight pending audit items without disturbing
% the body text style. Suppress via \tectsuppressmarks (define empty).
\providecommand{\tectauditflag}[1]{\textcolor{red}{\textbf{[AUDIT-FLAG: #1]}}}
\providecommand{\tecttodo}[1]{\textcolor{orange}{\textbf{[TODO: #1]}}}
\providecommand{\tecthonestscope}[1]{\textit{Honest scope: #1.}}

% =====================================================================
% End of TECT canonical preamble
% =====================================================================


\begin{document}

\title{An idea / mechanism note on per-generation chiral consistency
from combined Ward, Einstein--Hilbert, Witten, and mod-$2$ constraints
in TECT \\
\normalsize{\textit{(IDEA / MECHANISM NOTE; NOT a current theorem
paper. External use FORBIDDEN at the present canonical tier per
Math314-AddC audit. Full theorem-level rewrite tracked under
\texttt{Q-2026-05-02-Pillar7-Three-Generation-Rewrite}.)}}}

\author{Jusang Lee}
\email{jtkor@outlook.com}
\affiliation{Independent researcher, Republic of Korea}

\date{\today}

\begin{abstract}
This document is an idea / mechanism note outlining a candidate
framework for per-generation chiral consistency in the Topological
Energy Condensate Theory (TECT) combining four independent
consistency audits: (1) global Ward identities at each generation
(Math47); (2) preservation of those Ward currents under nonlinear
coupling to a fluctuating background metric (Math48); (3) the Witten
$\mathrm{SU}(2)_L$ global anomaly mod-$2$ constraint
(Math49b-v3, after \cite{Witten1982}); and (4) the mod-$2$
spectral-flow constraint on the chiral-fermion Dirac operator
(Math49c-v3). \emph{Honest scope and audit-flag (Math314-AddC).} The
present note does \emph{not} establish that the four conditions
together force the SM three-generation count: each is a per-generation
consistency \emph{requirement}, not a uniqueness theorem. The
generation-count question --- ``why exactly three?''  --- remains the
canonical Pillar~7 \textbf{T1 \textsc{open}} problem in TECT. The
correction history is recorded in the file header (Rev~4, 2026-05-06).
The companion paper~Paper-07-ext records the symbolic-exact
cancellation of the five SM perturbative gauge anomalies and the
$\mathrm{Tr}(\mathrm{U}(1)_Y)$ gravitational mixed anomaly per
generation; that result is independent of $N_g$ and is the canonical
SO(10) anomaly-cancellation audit of the TECT Pillar~4 embedding.
External use of this note is FORBIDDEN at the present canonical tier;
a full theorem-level rewrite of any three-generation-forcing claim is
tracked under
\texttt{Q-2026-05-02-Pillar7-Three-Generation-Rewrite}.
\end{abstract}

\maketitle

% =====================================================================
\section{Introduction}\label{sec:intro}
% =====================================================================
The existence of exactly three SM fermion generations is one of the
deepest empirical puzzles of the Standard Model. No first-principles
derivation has been universally accepted; the count $N_g = 3$ remains a
phenomenological input. TECT formulates the question in a constructive
form --- can the topological-condensate vacuum be required, by some
combination of field-theoretic consistency conditions, to admit only
$N_g = 3$? --- and answers it, at the present canonical tier, in the
\emph{negative}: the per-generation consistency audits identified here
are necessary, but they are not collectively known to be sufficient.

The present mechanism note assembles the four audit conditions that, in
combination, are the natural starting point for any future
three-generation-forcing argument inside TECT. They are: the global
Ward identity per generation (\S\ref{sec:ward}), the preservation of
that Ward current under nonlinear coupling to a fluctuating background
metric (\S\ref{sec:eh}), the Witten $\mathrm{SU}(2)_L$ global anomaly
mod-$2$ constraint (\S\ref{sec:witten}), and the mod-$2$ spectral-flow
constraint on the chiral-fermion Dirac operator (\S\ref{sec:mod2}).
None of them, individually or jointly, is known to force $N_g = 3$ at
the level of the present construction.

The paper is therefore positioned as an idea / mechanism note on
\emph{which audits} a future theorem-level construction must pass; it
is not, and does not claim to be, the theorem-level construction
itself.

% =====================================================================
\section{Setting and main claim (recast as conjecture)}\label{sec:setup}
% =====================================================================
For each generation $g$ of the SM, define the chiral fermion content
in the standard $\mathrm{SU}(3)_c\times \mathrm{SU}(2)_L\times
\mathrm{U}(1)_Y$ representation:
\begin{equation}\label{eq:sm-content}
  Q_L^g \in (\mathbf{3},\mathbf{2})_{1/6},\quad
  u_R^{c,g}\in (\overline{\mathbf{3}},\mathbf{1})_{-2/3},\quad
  d_R^{c,g}\in (\overline{\mathbf{3}},\mathbf{1})_{1/3},\quad
  L_L^g\in (\mathbf{1},\mathbf{2})_{-1/2},\quad
  e_R^{c,g}\in (\mathbf{1},\mathbf{1})_{1},
\end{equation}
together with the right-handed neutrino singlet
$\nu_R^{c,g}\in (\mathbf{1},\mathbf{1})_0$ (the RHN; SM hypercharge
$Y=0$) when SO(10) is the embedding GUT. The global Ward current of
generation $g$ for any conserved $\mathrm{U}(1)$ flavour symmetry $X$
is $j^{\mu}_{X,g}=\overline{\Psi}^g\gamma^{\mu}X\Psi^g$, satisfying
$\partial_\mu j^{\mu}_{X,g}=0$ at the classical level; quantum
corrections modify this only via gauge anomalies, which are audited in
the companion paper Paper-07-ext.

\begin{conjecture}[Per-generation quantum-consistency four-condition
audit; \textbf{T2 \textsc{conjecture}}]\label{conj:main}
For each generation $g\in\{1,\dots,N_g\}$ of the SM matter content
\eqref{eq:sm-content}, the following four conditions are
\emph{necessary} for quantum consistency of the per-generation chiral
sector:
\begin{itemize}
  \item[(\textbf{Ward-$g$})] global Ward identity per Math47 (see
    \S\ref{sec:ward}): $\partial_\mu j^{\mu}_{X,g} = 0$ for every
    classically conserved global $\mathrm{U}(1)_X$ symmetry, modulo
    the gauge anomalies that cancel between generations as audited in
    Paper-07-ext;
  \item[(\textbf{EH-$g$})] preservation of the Ward current
    $j^{\mu}_{X,g}$ under nonlinear coupling to a fluctuating
    background metric, per the one-loop Math48 audit
    (\S\ref{sec:eh});
  \item[(\textbf{Witten-tot})] the \emph{total} chiral
    $\mathrm{SU}(2)_L$ doublet count of the theory is even, per the
    Witten $\mathrm{SU}(2)$ global anomaly (Math49b-v3 +
    \cite{Witten1982}; \S\ref{sec:witten});
  \item[(\textbf{Mod-2-$g$})] the mod-$2$ spectral-flow index of the
    chiral Dirac operator on the BCC shell, per the Math49c-v3
    audit (\S\ref{sec:mod2}), is even for each generation.
\end{itemize}
The four-condition conjunction
$(\text{Ward-}g)\wedge(\text{EH-}g)\wedge
(\text{Witten-tot})\wedge(\text{Mod-2-}g)$ is necessary for the
per-generation chiral sector to admit a well-defined quantum
extension.

\textbf{Honest scope.} The conjunction is \emph{not} known to be
\emph{sufficient} to force the value of $N_g$. The generation-count
question --- ``does the four-condition conjunction admit solutions
only at $N_g = 3$?'' --- is the canonical Pillar~7 \textbf{T1
\textsc{open}} problem and is the subject of the rewrite obligation
\texttt{Q-2026-05-02-Pillar7-Three-Generation-Rewrite}.
\end{conjecture}

\begin{remark}[On the difference from the earlier draft]\label{rmk:diff}
Earlier drafts of this paper asserted that the conjunction
\emph{forces} $N_g = 3$ at \textbf{T7 \textsc{proved}}. That assertion
is withdrawn in Rev~4 because (i) the Witten $\mathrm{SU}(2)$ global
anomaly is even for $N_g\in\{1,2,3,\dots\}$ (the per-generation
doublet count is $4$, always even), and therefore does not
discriminate $N_g = 3$; (ii) the formula
$A_{\mathrm{grav}}=\sum_f Q_f^3$ used in the earlier ``Step 2'' is
\emph{not} the SM gravitational anomaly coefficient (the actual mixed
gravitational anomaly is $\mathrm{Tr}(\mathrm{U}(1)_Y)$, with one
power of $Y$, which vanishes per generation independent of $N_g$;
this is the canonical SO(10) audit of Paper-07-ext); (iii) the
mod-$2$ parity expression $(-1)^{N_g(N_g-1)/2}$ in the earlier
``Step 4'' is non-standard and not tied to any canonical TECT
construction. The corrected statement is recorded in
Conjecture~\ref{conj:main}.
\end{remark}

% =====================================================================
\section{Audit~1 --- per-generation Ward identity (Math47)}
\label{sec:ward}
% =====================================================================
Math47 establishes that, for the SM coupled to the TECT background
condensate, the global Ward identity for any classically conserved
$\mathrm{U}(1)_X$ flavour current of generation $g$ is preserved at
all loop orders, provided the gauge-anomaly cancellation conditions
of Paper-07-ext are satisfied for the full chiral spectrum. The
precise theorem is the per-generation independence statement: the
Ward current $j^{\mu}_{X,g}$ does not mix with $j^{\mu}_{X,g'}$ for
$g\neq g'$ at tree level, and the loop-level mixing is governed by
the same gauge-anomaly trace formulas audited in Paper-07-ext.

\begin{remark}[What Math47 does \emph{not} establish]
\label{rmk:math47-limits}
Math47 does not by itself force any particular value of $N_g$: the
per-generation Ward identity holds at every $g$ \emph{individually},
and the gauge-anomaly cancellation it depends on is, per
the six-coefficient table (eq.~6 of Paper-07-ext), $N_g$-independent.
\end{remark}

% =====================================================================
\section{Audit~2 --- preservation under nonlinear EH coupling
(Math48)}\label{sec:eh}
% =====================================================================
When the SM matter is coupled to a fluctuating background metric in
the linearised regime $g_{\mu\nu} = \eta_{\mu\nu} + h_{\mu\nu}$, the
mixed gravitational anomaly potentially obstructs the conservation of
each Ward current. Math48 establishes, at one-loop, the nonlinear
preservation of $\partial_\mu j^{\mu}_{X,g} = 0$ provided the
gravitational mixed anomaly $\mathrm{Tr}(\mathrm{U}(1)_X)$ vanishes
for every gauged $X$. For $X = Y$ (the SM hypercharge), this
condition is the audit-target of Paper-07-ext, where it is proven
symbolically exact zero per generation
\eqref{eq:trY-zero-7ext-ref}; for $X = B-L$ (the SO(10) extension),
it is closed by the RHN.

\begin{remark}[Removal of the non-standard $A_{\mathrm{grav}} = \sum_f
Q_f^3$ formula]\label{rmk:remove-Q3}
The earlier draft (Rev~3) introduced the formula
$A_{\mathrm{grav}} = \sum_f Q_f^3$ (sum over fermions of the cube of
the electric charge $Q$) as the ``gravitational anomaly coefficient''.
This is incorrect: the gravitational mixed anomaly in four dimensions
is the trace $\mathrm{Tr}(\mathrm{U}(1)_X)$ for the gauged Abelian
charge $X$ (one power, not three), arising from the
graviton--graviton--gauge-boson triangle. The pure-gauge anomaly
$\mathrm{Tr}(\mathrm{U}(1)_X^3)$ (three powers) is a
\emph{separate} coefficient, and for $X = Q$ (electric charge) is the
non-canonical combination of $\mathrm{U}(1)_Y$ and SU(2)$_L$
contributions which is normally evaluated only in the broken
electroweak phase below the Higgs vev. The earlier draft conflated
these two coefficients. In Rev~4 the formula is removed and the audit
is the standard $\mathrm{Tr}(\mathrm{U}(1)_Y)$ as recorded in
Paper-07-ext, equation~\eqref{eq:trY-zero-7ext-ref}.
\end{remark}

A local stub of the canonical identity, proved as
in the six-coefficient table of the companion paper~Paper-07-ext, is
recorded here for self-contained citation:
\begin{equation}\label{eq:trY-zero-7ext-ref}
  \mathrm{Tr}(\mathrm{U}(1)_Y) \;=\; 0 \quad\text{per generation}
  \qquad\text{[symbolic identity, proven in Paper-07-ext]}.
\end{equation}

% =====================================================================
\section{Audit~3 --- Witten $\mathrm{SU}(2)_L$ global anomaly
(Math49b-v3, after Witten 1982)}\label{sec:witten}
% =====================================================================
Witten \cite{Witten1982} proved that a four-dimensional chiral gauge
theory with gauge group $\mathrm{SU}(2)$ is consistent at the
non-perturbative level only if the total number of chiral
$\mathrm{SU}(2)$ doublets is even. The argument uses
$\pi_4(\mathrm{SU}(2)) = \mathbb{Z}_2$: a homotopically nontrivial
$\mathrm{SU}(2)$ gauge transformation generates a sign in the
fermionic path-integral measure that vanishes only if the total
doublet count is even.

For the SM at $N_g$ generations, the chiral $\mathrm{SU}(2)_L$
doublets per generation are: the three quark doublets (one per
colour) of $Q_L$, plus the lepton doublet $L_L$. The total per-generation
doublet count is therefore $3+1 = 4$, which is even. The total over
$N_g$ generations is $4 N_g$, which is even for every positive
integer $N_g$.

\begin{remark}[Witten does not discriminate $N_g = 3$]
\label{rmk:witten-noselect}
The Witten $\mathrm{SU}(2)_L$ global anomaly is satisfied for every
$N_g\geq 1$. It does not discriminate $N_g = 1$ from $N_g = 3$. The
earlier draft assertion that ``the coefficient is zero iff
$N_g = 3$'' is corrected in Rev~4: the coefficient (the
$\mathbb{Z}_2$ class) is zero for every $N_g$.
\end{remark}

% =====================================================================
\section{Audit~4 --- mod-$2$ spectral flow on the BCC shell
(Math49c-v3)}\label{sec:mod2}
% =====================================================================
Math49c-v3 examines the chiral Dirac operator on the BCC shell of the
TECT condensate vacuum. Its integer-valued Atiyah--Singer index, per
generation, is given by a topological-charge sum on the shell;
quantum consistency of the chiral path integral requires the
mod-$2$ reduction of this index to be even (so that the fermionic
Pfaffian is well-defined as a real-valued square root of the chiral
Dirac determinant).

\begin{remark}[Mod-$2$ does not, in the present construction, force
$N_g = 3$]\label{rmk:mod2-noselect}
The mod-$2$ spectral-flow index of Math49c-v3 is, in the present
canonical construction, an integer that is even for the SM SU(5)
content of $\overline{\mathbf{5}}\oplus \mathbf{10}$ per generation;
this is consistent with $N_g = 1, 2, 3, \dots$\ alike. A construction
in which the index parity is sensitive to $N_g$ is conceivable but is
not part of the present Math49c-v3 archive; if completed, it would
be the central technical step of any future
three-generation-forcing theorem.
\end{remark}

% =====================================================================
\section{Closure and honest verdict}\label{sec:closure}
% =====================================================================
The four audits above are individually closed at \textbf{T6
\textsc{proved conditional}} or stronger (Ward, EH, Witten, mod-$2$
each rest on textbook or Math47/Math48/Math49b-v3/Math49c-v3
canonical archive). \emph{Their conjunction is necessary} for the
per-generation chiral sector to admit a well-defined quantum
extension.

\textbf{The conjunction is not known to be sufficient to force the SM
generation count.} Specifically:
\begin{itemize}
  \item Witten $\mathrm{SU}(2)_L$ is satisfied for every $N_g\geq 1$;
  \item Mod-$2$ spectral flow, in the present Math49c-v3 construction,
    is satisfied for every $N_g$ for which the chiral SU(5) content
    sums to an even index --- which is the case for every $N_g$;
  \item Ward and EH are per-generation conditions that hold
    individually for each $g\in\{1,\dots,N_g\}$ and are
    $N_g$-independent;
  \item the perturbative SM gauge anomalies and
    $\mathrm{Tr}(\mathrm{U}(1)_Y)$ gravitational anomaly cancel
    \emph{per generation} (Paper-07-ext, equation
    \eqref{eq:trY-zero-7ext-ref} and the six-coefficient table of
    that paper), again $N_g$-independent.
\end{itemize}

The honest verdict is therefore that the four-condition consistency
audit, combined with the Paper-07-ext anomaly cancellation, fixes
$N_g$ to be a positive integer (each generation must satisfy the
audit) but \emph{does not} fix $N_g = 3$. The phenomenological value
$N_g = 3$ remains the empirical input, and the Pillar~7 question
``why exactly three?''  is canonical \textbf{T1 \textsc{open}} in
TECT, tracked under
\texttt{Q-2026-05-02-Pillar7-Three-Generation-Rewrite}.

% =====================================================================
\section{Discussion and outlook}\label{sec:discussion}
% =====================================================================
The mechanism note recorded here serves three purposes within the
TECT Pillar~7 programme:
\begin{enumerate}
  \item to enumerate the necessary per-generation consistency audits
    that any future three-generation-forcing theorem must pass;
  \item to record the Math314-AddC corrections to the earlier
    over-claims (Rev~3 had asserted $N_g = 3$ forced at \textbf{T7
    \textsc{proved}}, an over-claim now retracted);
  \item to provide a clean target for the
    \texttt{Q-2026-05-02-Pillar7-Three-Generation-Rewrite} obligation,
    which would supply a fifth condition (a TECT-internal
    construction sensitive to $N_g$) sufficient to force $N_g = 3$
    in conjunction with the four conditions above.
\end{enumerate}
The fifth condition, if ever constructed, would most likely arise
from the BCC shell topology --- the integer-valued index of the
Math49c-v3 chiral Dirac operator restricted to the three-band shell
spectrum (Math234), whose three-fold degeneracy might be tied to a
TECT-internal Chern number or flux quantum. Such a construction is a
canonical Pillar~7 T1 OPEN item and is not part of the present note.

% =====================================================================
\begin{thebibliography}{99}
\bibitem{Witten1982} E.~Witten, ``An $\mathrm{SU}(2)$ anomaly,''
  Phys.\ Lett.\ B \textbf{117}, 324 (1982).
\bibitem{Math47} TECT internal note Math47,
  ``Per-generation Ward identity in the SM coupled to TECT
  condensate,'' TECT canonical archive.
\bibitem{Math48} TECT internal note Math48,
  ``Nonlinear EH coupling and Ward-current preservation at one-loop,''
  TECT canonical archive.
\bibitem{Math49bv3} TECT internal note Math49b-v3,
  ``Witten $\mathrm{SU}(2)_L$ global anomaly audit for the TECT
  chiral spectrum,'' TECT canonical archive.
\bibitem{Math49cv3} TECT internal note Math49c-v3,
  ``Mod-$2$ spectral flow on the BCC shell chiral Dirac operator,''
  TECT canonical archive.
\bibitem{Math314AddC} TECT internal note Math314-AddC,
  ``Hostile-referee audit of Wave-4/5 papers Paper-06..08; rewrite
  obligation Q-2026-05-02-Pillar7-Three-Generation-Rewrite,''
  TECT canonical archive.
\end{thebibliography}

\begin{acknowledgments}
The author thanks the Math314-AddC hostile-referee audit (2026-05-02)
for catching the Rev~3 over-claims corrected in Rev~4. This work is
part of TECT, canonical archive at
\url{https://github.com/TECT-OpenPhysics/TECT}.
\end{acknowledgments}

\end{document}
